\documentclass[a4paper,12pt]{article}

\usepackage{mystyle}

\usepackage{gensymb}
\usepackage{scalerel}
\usepackage{stackengine}

% \usepackage{skull}  % skull
\usepackage{halloweenmath}  % \bigpumpkin, skull (https://tug.ctan.org/info/symbols/comprehensive/symbols-a4.pdf -- Table 76)

\usepackage{tikzsymbols}

% https://tex.stackexchange.com/questions/3266/how-do-i-use-a-circle-as-a-math-accent-larger-than-mathring
% https://tex.stackexchange.com/a/3270/135045
\usepackage{accents}

\renewcommand{\mathring}[1]{\accentset{\circ}{#1}}


\graphicspath{ {images/} }


% https://tex.stackexchange.com/questions/5461/is-it-possible-to-change-the-size-of-an-arrowhead-in-tikz-pgf
\usetikzlibrary{arrows.meta}


\DeclareMathOperator{\Image}{Im}

\definecolor{pink}{RGB}{218, 3, 174}
\definecolor{violet}{RGB}{148, 0, 211}
\definecolor{green}{RGB}{0, 153, 0}
\definecolor{orange}{RGB}{255, 153, 0}
\definecolor{blue}{RGB}{5, 73, 255}
\definecolor{cyan}{RGB}{31, 206, 203}
\definecolor{cyan2}{RGB}{0, 166, 147}
\definecolor{cyangreen}{RGB}{0, 155, 118}
\definecolor{cyangreen2}{RGB}{0, 109, 91}


% https://tex.stackexchange.com/a/101138/135045

\newcommand\widesim[1]{\ThisStyle{%
  \setbox0=\hbox{$\SavedStyle#1$}%
  \stackengine{-.1\LMpt}{$\SavedStyle#1$}{%
    \stretchto{\scaleto{\SavedStyle\mkern.2mu\sim}{.5150\wd0}}{.6\ht0}%
  }{O}{c}{F}{T}{S}%
}}


\newcommand{\BigMiddleThree}{\;\left|\vphantom{\begin{pmatrix} 0\\0\\0 \end{pmatrix}}\right.\;}
\newcommand{\BigMiddleFour}{\;\left|\vphantom{\begin{pmatrix} 0\\0\\0\\0 \end{pmatrix}}\right.\;}


% https://tex.stackexchange.com/questions/63531/how-to-write-quotation-marks-in-math-environment
\DeclareMathSymbol{\mlq}{\mathord}{operators}{``}
\DeclareMathSymbol{\mrq}{\mathord}{operators}{`'}


\DeclareMathOperator{\Imag}{Im}


% https://tex.stackexchange.com/questions/544453/undefined-control-sequence-after-paragraph
\renewcommand{\paragraph}[1]{\noindent\textbf{#1}\quad}


% https://tex.stackexchange.com/questions/36851/skipping-line-after-proof-in-proof-environment#comment73553_36851
\newcommand{\proofindent}{\hspace*{\fill}\par\vspace{0.5em}\noindent}


% https://tex.stackexchange.com/questions/4813/extendible-equals-sign
\makeatletter
\newcommand*{\Relbarfill@}{\arrowfill@\Relbar\Relbar\Relbar}
\newcommand*{\xeq}[2][]{\ext@arrow 0055\Relbarfill@{#1}{#2}}
\makeatother


% https://tex.stackexchange.com/questions/279100/typeset-the-shrug-%C2%AF-%E3%83%84-%C2%AF-emoji
\newcommand{\shrug}[1][]{%
\begin{tikzpicture}[baseline,x=0.8\ht\strutbox,y=0.8\ht\strutbox,line width=0.125ex,#1]
  \def\arm{(-2.5,0.95) to (-2,0.95) (-1.9,1) to (-1.5,0) (-1.35,0) to (-0.8,0)};
  \draw \arm;
  \draw[xscale=-1] \arm;
  \def\headpart{(0.6,0) arc[start angle=-40, end angle=40,x radius=0.6,y radius=0.8]};
  \draw \headpart;
  \draw[xscale=-1] \headpart;
  \def\eye{(-0.075,0.15) .. controls (0.02,0) .. (0.075,-0.15)};
  \draw[shift={(-0.3,0.8)}] \eye;
  \draw[shift={(0,0.85)}] \eye;
  % draw mouth
  \draw (-0.1,0.2) to [out=15,in=-100] (0.4,0.95); 
\end{tikzpicture}}



% https://tex.stackexchange.com/a/314638/135045
% \newcommand{\diff}{\mathop{}\!d\!}
\newcommand{\diff}{\mathop{}\!d}


\author{Алексеев Василий}


\title{Семинар 10}
\date{18 ноября 2024}


\begin{document}
  \maketitle
  
  \tableofcontents

  \thispagestyle{empty}
  
  \newpage
  
  
  
  \vspace*{\fill}
  
  \noindent
  \emph{
    К формулировкам и доказательствам (если такие вообще приводятся) стоит относиться критически.
    Основное в этом конспекте~---~решение задач (но ``критичность'' и здесь лучше не отключать).
    За строгой, ясной и последовательной теорией лучше обращаться к ``нормальным'' источникам.
    (Например, к лекциям.)
  }
  
  \vspace*{\fill}
  
  \thispagestyle{empty}
  
  \newpage
  
  
  \pagenumbering{arabic}


  \section{Равномерная непрерывность}
  
  %\begin{figure}[ht]
    %\centering
    %\includegraphics[width=0.6\linewidth]{images/Apples}
    %
    %\caption{
      %Одно яблоко, два яблока, ...~---~натуральные числа используются при счёте предметов.
    %}
    %\label{fig:naturals}
  %\end{figure}
  
  
  % TODO:
  % [ ] Определение
  % [ ] Примеры (x, x^2, sin x)
  % [ ] Непрервность на отрезке -> равномерная непрерывность
  % [ ] (?) Модуль непрерывности
  
  
  \subsection{С1, \S 12, \textnumero 17}
  \label{p:1-12-17}
  
  Пусть функция~$f$ равномерно непрерывна на отрезках~$[a, b]$ и~$[b, c]$.
  Доказать, что в таком случае она будет равномерно непрерывна и на отрезке~$[a, c]$.
  
  \begin{solution}
    % TODO: pic
    Надо доказать, что  % TODO: ref
    \[
      \forall \eps > 0\ \exists \delta > 0\colon \forall x, x'\colon |x - x'| < \delta \to |f(x) - f(x')| < \eps
    \]

    Если $x,\ x'$ оба лежат либо на $[a, b]$, либо на $[b, c]$, то всё работает (на каждом из отрезков существует подходящий~$\delta$~---~пусть это $\delta_l$ и $\delta_r$ соответственно).

    Рассмотрим случай, когда, например, $x \hm\in [a, b]$, а $x' \hm\in [b, c]$.
    Считаем, что $|x \hm- x'| \hm< \delta$ (разность можно выбрать такой маленькой, как хотим).
    Посмотрим на разность значений функции в этих точках:
    \[
      |f(x) - f(x')| = \Bigl|\bigl(f(x) - f(b)\bigr) + \bigl(f(b) - f(x')\bigr)\Bigr|
        \leq |f(x) - f(b)| + |f(b) - f(x')| < 2 \eps
    \]
    где в последнем неравенстве считали, что $\delta \hm< \min{(\delta_l, \delta_r)}$, где $\delta_l \hm= \delta_l(\eps)$ и $\delta_r \hm= \delta_r(\eps)$.

    Итого, при $\delta \hm< \min{\bigl(\delta_l(\eps/2), \delta_r(\eps/2)\bigr)}$ получаем ограниченность разности значений функций $|f(x) \hm- f(x')| \hm< \eps$ для всех пар~$x,\ x'$, таких что $|x \hm- x'| \hm< \delta$.
  \end{solution}
  
  
  
  \subsection{С1, \S 12, \textnumero 23}
  
  Функция~$f$ непрерывна на~$[a, +\infty)$ и существует конечный $\lim\limits_{x \to +\infty} f(x)$.
  Доказать, что~$f$ равномерно непрерывна на~$[a, +\infty)$.
  
  \begin{solution}
    \emph{Идея}.
    Существует конечный предел на бесконечности~---~значит, по мере увеличения~$x$ значения~$f(x)$ становятся всё ближе к этому предельному значению, а потому и ближе друг к другу.  % TODO: pic
    Всегда можно подобрать такую окрестность бесконечности, чтобы разность $\bigl|f(x) \hm- f(x')\bigr|$ для точек~$x,\ x'$ из этой окрестности была сколь угодно маленькой...
    А на оставшемся отрезке функция будет непрерывна, а потому и равномерно непрерывна.

    Итак, раз существует конечный предел~$f(x)$ на бесконечности (равный для определённости~$A$), значит:
    \[
      \forall \eps > 0\ \exists \delta_r > 0\colon \forall x > b \equiv \frac{1}{\delta_r} \to \bigl|f(x) - A\bigr| \hm< \eps
    \]

    На отрезке же~$[a, b]$ функция~$f$ непрерывна.
    Поэтому на нём она будет и равномерно непрерывна.  % TODO: ref
    То есть найдётся $\delta_l$, такой что разность $\Delta f$ для любой пары точек отрезка, находящихся друг от друга на расстоянии менее~$\delta_l$, будет по модулю меньше~$\eps$.

    Аналогично задаче~(\ref{p:1-12-17}), можно рассмотреть случай, когда одна из точек~$x,\ x'$ находится на~$[a, b]$, а другая на~$[b, +\infty)$.
    При расстоянии между ними $\delta \hm= \min(\delta_l, \delta_r)$ разность $\Delta f$ для этих точек будет по модулю ограничена~$2 \eps$.

    Итого, при $\delta \hm= \min\bigl(\delta_l(\eps/2), \delta_r(\eps/2)\bigr)$ разность $\Delta f$ для любой пары точек~$x,\ x' \hm\in [a, +\infty)$, таких что $|x \hm- x'| \hm< \delta$, будет ограничена по модулю величиной~$\eps$.
  \end{solution}



  % \subsection{С1, \S 12, \textnumero 25}
  
  % Доказать, что для равномерной непрерывности функции~$f$ на ограниченном интервале~$(a, b)$ необходимо и достаточно, чтобы функция была непрерывна на~$(a, b)$ и чтобы существовали конечные пределы $\lim\limits_{x \to a + 0} f(x)$ и $\lim\limits_{x \to b - 0} f(x)$.
  % % TODO: про "конечные" в условии не было
  
  % \begin{solution}
  
  % \end{solution}
  
  
  
  \subsection{С1, \S 12, Т1(а)}
  
  Пусть функция~$f$ дифференцируема на~$[a, +\infty)$ и её производная~$f'$ ограничена.
  Доказать, что~$f$ равномерно непрерывна на~$[a, +\infty)$.
  
  \begin{solution}
    \emph{Идея}.
    Если производная в любой точке~$x$ ограничена, то разность $f(x') \hm- f(x)$ будет стремиться к нулю при стремлении $x'$ к~$x$:
    \[
      f'(x) = \lim_{x' \to x} \frac{f(x') - f(x)}{x' - x} < C
    \]
  \end{solution}
  
  
  
  \subsection{С1, \S 12, Т2(а, б)}
  
  Исследовать функции на равномерную непрерывность на луче~$(0, +\infty)$:
  \[
    \mbox{а})\ f(x) = x^3 \sin{\frac{1}{x}},\quad \mbox{б})\ f(x) = x e^{\sin x}
  \]
  
  \begin{solution}
  
  \end{solution}
  

\end{document}
